\section{Factor Structure: The Volatility-Gated PCA Adapter}
\label{sec:pca_factors}

% Arc position: immediately after the two-block (backbone + exogenous) ridge
% of Section~\ref{sec:dense_weak}. The PCA factor block is the third rung of
% the block ladder. This section subsumes the factor construction formerly in
% the transmission section (the lead--lag / operator machinery is dropped)
% and adds the P-model mechanism battery of 2026-08-09/10. All mechanism arms
% are complete 91-chunk runs on the frozen panel (chunks 9-99, 248,686 joined
% OOS rows), harvested through score_unification.py; paired Diebold-Mariano on
% common rows. Diagnostics: results/p_hypothesis_analysis.json,
% results/frame_alignment_diagnostics.csv. Arm definitions in
% src/unification.py (blk3_*Spike* family). Writeup handoff:
% writeup/P_MODEL_RESULTS_BRIEF.md. 2026-08-10.

The two blocks assembled so far read the current bar: the target's own
history and the exogenous panel. This section adds a third block built from
the panel's eigenstructure, in four steps, and then takes the block apart,
in three steps, to see where its value comes from. The construction is one
pass of principal-component analysis followed by one scaling step. The
decomposition shows that the block is a \emph{volatility-gated PCA
adapter}: a low-rank adaptation of the exogenous loadings along the
panel's principal directions, gated per direction by trailing realized
volatility.

\subsection{The Construction, in Four Steps}
\label{sec:pca_frame}

\paragraph{Step 1: freeze a frame of principal directions.}
Take the panel's base-feature vector $x(t)$ at bar $t$ and keep only its
live columns --- the set $J$ of directions whose sample standard deviation
exceeds the degeneracy threshold over the first training window,
\[
\mathcal{T}_0 \;=\; \{\, W+1, \dots, 2W \,\},
\qquad W = 24000 \text{ bars}.
\]
Standardize each live column by its frame-window mean $\mu_j$ and standard
deviation $\sigma_j$, form the correlation matrix of the live columns over
$\mathcal{T}_0$, and take its eigendecomposition,
\[
R\, v_k \;=\; \lambda_k v_k,
\qquad \lambda_1 \ge \lambda_2 \ge \dots \ge \lambda_p .
\]
The frame is the top-$q$ eigenvectors $v_1, \dots, v_q$. Every quantity in
this step --- $\mu_j$, $\sigma_j$, $v_i$ --- is estimated once, from
feature data alone, and the frame is then \textbf{frozen permanently}
before any forecasting begins: a fixed coordinate system, not a current
one. The width $q = 40$ is a representative of a plateau, not a tuned
optimum. The first-window spectrum's participation ratio puts the
effective span at $\approx 20$ dimensions; the raw Marchenko--Pastur edge
leaves roughly 22 live eigenvalues above the independent-noise level; a
dependence-adjusted edge, deflating the sample by the factors' integrated
autocorrelation time (median $175$ bars), leaves only 5 eigenvectors
individually identified. The data pin the subspace the mid-spectrum
eigenvectors span, not the basis inside it, and the width experiments
confirm it: performance depends on covering the span, not on where inside
it the truncation falls.

\paragraph{Step 2: project the panel onto the frame.}
At every bar $t$, the factor score $G_i(t)$ is the projection of the
bar's feature vector onto frame direction $i$, re-standardized by the
projection's own frame-window moments $m_i$ and $s_i$:
\[
G_i(t) \;=\; \frac{1}{s_i} \Bigl( \sum_{j \in J} v_{i,j}\, z_j(t) - m_i \Bigr),
\qquad i = 1, \dots, q,
\]
where $z_j(t)$ is the standardized feature from Step 1. Because the frame
is frozen, each score is a \emph{fixed} linear combination of columns
already in the design --- this linearity is the whole mechanism, as
Section~\ref{sec:pca_mechanism} shows.

\paragraph{Step 3: gate each direction by its trailing volatility.}
Each factor score is scaled by its own trailing standard deviation,
computed causally over a rolling window that uses only past bars: the
column entering the design at bar $t$ is $s_i(t)\, G_i(t)$, with $s_i(t)$
the realized volatility of direction $i$'s score up to $t$. Directions
that have been volatile recently are amplified; quiet directions are
suppressed. This step looks like housekeeping. It is not:
Section~\ref{sec:pca_mechanism} shows it carries the majority of the
block's value.

\paragraph{Step 4: shrink the scores as a block of their own.}
The $q$ gated scores enter the ridge as a third block alongside the
backbone and the exogenous panel, with its own penalty, selected by the
same per-block causal tuning as every other block. As in the previous
sections, per-block penalties are implemented by column scaling, so one
solver call fits all three blocks.

\subsection{What the Block Is, in Three Steps}
\label{sec:pca_mechanism}

\paragraph{Step 1: the block adds no information.}
Because Step 2 is a fixed linear map, every factor column lies exactly in
the span of the backbone and exogenous columns. Projecting the block's
columns onto that span leaves a residual squared-norm share of
$3.6 \times 10^{-30}$ --- machine roundoff in float64. The block cannot
add information; it can only change the estimator's geometry. The right
question is therefore not ``what do the factors know?'' but ``what does
the block do to the fit?''

\paragraph{Step 2: without the gate, the block is exactly a spiked prior.}
Freeze the gate ($s_i(t) \equiv s_i$, constant) and write the model with
the factor block explicit:
\[
y \;=\; a + B\beta_B + Z\beta_Z + G\gamma + \varepsilon,
\qquad G = ZA,
\]
where $A$ is the frozen frame map from Steps 1--2 and the ridge penalties
are $\lambda_Z \lVert \beta_Z \rVert^2$ on the exogenous coefficients and
$\gamma^{\top} D_G\, \gamma$ on the factor coefficients. Now substitute
\[
\delta \;=\; \beta_Z + A\gamma .
\]
The two exogenous terms collapse into one, the two quadratic penalties
combine, and the fit is \emph{identical} --- same normal equations, same
forecasts, same QLIKE --- to the two-block generalized ridge
\[
y \;=\; a + B\beta_B + Z\delta + \varepsilon,
\qquad
\delta \sim N\!\bigl(0,\; P^{-1}\bigr),
\qquad
P^{-1} \;=\; \lambda_Z^{-1} I + A D_G^{-1} A^{\top}.
\]
Read as a prior, this says the exogenous coefficient vector is not
shrunk uniformly: it carries a base isotropic prior plus a rank-$q$
covariance spike along the frame's principal directions. That is exactly
the low-rank-update form of a LoRA-type adapter --- base weights plus a
rank-$q$ correction --- with the adapter basis fixed to the PCA
eigenvectors instead of learned. There is nothing to fit beyond this
identity: the three-block frozen regression \emph{is} the $P$-model,
reparametrized.

\paragraph{Step 3: the gate is the part no static prior can express.}
Restore Step 3. The gated scores are
\[
G \;=\; D_s\, Z A,
\]
with $D_s$ a diagonal matrix of time-varying scale factors. There is no
constant map $A'$ with $D_s Z A = Z A'$, so no time-invariant Gaussian
prior on $\delta$ --- no choice of $P$ --- can reproduce the gated block.
What the gate adds is a time-varying coefficient path,
\[
\delta_t \;=\; \beta_Z + A\,\mathrm{diag}(s_t)\,\gamma :
\]
a rank-$q$ adaptation of the exogenous loadings whose per-direction gain
$s_t$ is the trailing realized volatility of that principal direction.
Locally, direction $i$ at time $t$ enters the normal equations with
weight $s_{t,i}^2$, so its effective shrinkage
\[
\kappa_i(t) \;=\; \frac{s_{t,i}^2}{s_{t,i}^2 + d_i / W}
\]
relaxes exactly when the direction has been volatile, and tightens when
it has been quiet. The static $P$ of Step 2 is the constant-gate special
case.

The measured values quantify the split. On the frozen panel the best
static-spike arm --- the full $P$-model, tuned --- scores QLIKE
$0.21990$ against the backbone-plus-exogenous ridge's $0.21997$: about a
fifth of the block's measured value. The gated production construction
reaches $0.21961$, and the frozen-to-gated increment on the production
frame is $-4.3 \times 10^{-4}$ at DM $t = -3.17$. The adapter's basis
and its gate are both necessary; the gate contributes the majority.

\subsection{Which PCA: The Frame Is the Mechanism}
\label{sec:pca_which_frame}

If the block is a prior spike along principal directions, one design
question remains: which variables the PCA of Step 1 is estimated on ---
the \emph{frame}. The answer is decisive, and it is the paper's main
evidence that the basis, not the machinery around it, carries the
signal.

First, a pure exogenous frame --- the same construction, the same tuned
grids, but with Step 1 restricted to the broad exogenous columns alone
--- delivers $\Delta\mathrm{QLIKE} = -6.8 \times 10^{-5}$ against the
plain backbone-plus-exogenous ridge, at DM $t = -0.76$: nothing. The
reason is geometric. Estimating ordinary rolling ridge coefficient
directions in five windows spread across the sample, and measuring how
much of each direction the top-$40$ subspace of each candidate frame
captures, the exogenous frame's share is $0.02$--$0.07$; the production
frame's subspace --- the mixed frame of HAR, session, and short-window
exogenous columns --- captures the dominant share of the coefficient
direction, and its eigenvectors load jointly on the two column families
(backbone/session loading mass $0.43$, exogenous $0.57$). The useful
principal directions are \emph{mixed-family} directions, and an
exogenous-only PCA does not contain them.

Three deletions sharpen the statement. Removing the cross-block
covariance of the induced joint prior --- splitting the spike into
separate backbone and exogenous parts, deleting the cross term
$A_B D_G^{-1} A_Z^{\top}$ --- costs $3.0 \times 10^{-5}$ at DM
$t = -2.48$: small, but the only statistically detectable effect among
the frozen-frame variants. Removing the HAR columns or the session
columns from the frame changes nothing (DM $t = +0.02$ and $-0.36$).
Permuting each column family independently within the frame window ---
destroying all cross-family time synchrony while preserving every
marginal distribution --- also changes nothing detectable (DM
$t = +1.15$). What matters is that Step 1's PCA sees both families'
cross-family correlation structure; which family members supply it, and
whether the two families are contemporaneously aligned within the
window, do not.

Finally, the frame and the gate interact rather than add. Trailing
scaling on the exogenous frame improves nothing detectable
($-1.1 \times 10^{-4}$, DM $t = -0.84$), and the production frame with a
frozen gate is worse than the exogenous frame with a frozen gate
($+1.4 \times 10^{-4}$, $t = +1.21$). Only the pair --- the right basis,
volatility-gated --- recovers the block's value.

\subsection{What Remains Outside the Adapter}
\label{sec:pca_outside}

Two components of the headline model are not part of this mechanism. The
frozen pairwise-product block of Section~\ref{sec:nonlinear_product}
adds $-7.5 \times 10^{-4}$ at DM $t = -1.76$ beyond the linear factor
construction --- roughly two-thirds of the final model's total gain over
the backbone-plus-exogenous ridge --- and is a genuinely nonlinear
channel, not a proxy for the adapter. And the exogenous block itself
remains where the model's forecast skill lives: the two-block ridge's
paired increment against the HAR-only benchmark, measured on common rows,
is $-5.8 \times 10^{-3}$ at DM $t = -12.0$ --- an order of magnitude
beyond every refinement this section reports.

\paragraph{Reading.}
The factor block's value decomposes into a basis and a gate. The basis
is a set of mixed-family principal directions of the panel (Step 1); the
gate is their trailing realized volatilities (Step 3). Neither works
alone; the static prior $P$ is the constant-gate special case; and once
the basis and the gate are in place there is no further machinery to
estimate. Two qualifications accompany every statistic in this section:
all inference is on a single out-of-sample path, and the mechanism arms
were selected among dozens of campaign arms, so the marginal
Diebold--Mariano values reported here (the cross-covariance deletion at
$t = -2.48$ in particular) should be read against that
multiple-comparison background. The frame-alignment contrast and the
frame-by-scaling interaction are the robust objects; the rest is
supporting detail.
